\IfFileExists{stacks-project.cls}{%
\documentclass{stacks-project}
}{%
\documentclass{amsart}
}

% For dealing with references we use the comment environment
\usepackage{verbatim}
\newenvironment{reference}{\comment}{\endcomment}
%\newenvironment{reference}{}{}
\newenvironment{slogan}{\comment}{\endcomment}
\newenvironment{history}{\comment}{\endcomment}

% For commutative diagrams we use Xy-pic
\usepackage[all]{xy}

% We use 2cell for 2-commutative diagrams.
\xyoption{2cell}
\UseAllTwocells

% We use multicol for the list of chapters between chapters
\usepackage{multicol}

% This is generally recommended for better output
\usepackage{lmodern}
\usepackage[T1]{fontenc}

% For cross-file-references
\usepackage{xr-hyper}

% Package for hypertext links:
\usepackage{hyperref}

% For any local file, say "hello.tex" you want to link to please
% use \externaldocument[hello-]{hello}
\externaldocument[introduction-]{introduction}
\externaldocument[conventions-]{conventions}
\externaldocument[sets-]{sets}
\externaldocument[categories-]{categories}
\externaldocument[topology-]{topology}
\externaldocument[sheaves-]{sheaves}
\externaldocument[sites-]{sites}
\externaldocument[stacks-]{stacks}
\externaldocument[fields-]{fields}
\externaldocument[algebra-]{algebra}
\externaldocument[brauer-]{brauer}
\externaldocument[homology-]{homology}
\externaldocument[derived-]{derived}
\externaldocument[simplicial-]{simplicial}
\externaldocument[more-algebra-]{more-algebra}
\externaldocument[smoothing-]{smoothing}
\externaldocument[modules-]{modules}
\externaldocument[sites-modules-]{sites-modules}
\externaldocument[injectives-]{injectives}
\externaldocument[cohomology-]{cohomology}
\externaldocument[sites-cohomology-]{sites-cohomology}
\externaldocument[dga-]{dga}
\externaldocument[dpa-]{dpa}
\externaldocument[sdga-]{sdga}
\externaldocument[hypercovering-]{hypercovering}
\externaldocument[schemes-]{schemes}
\externaldocument[constructions-]{constructions}
\externaldocument[properties-]{properties}
\externaldocument[morphisms-]{morphisms}
\externaldocument[coherent-]{coherent}
\externaldocument[divisors-]{divisors}
\externaldocument[limits-]{limits}
\externaldocument[varieties-]{varieties}
\externaldocument[topologies-]{topologies}
\externaldocument[descent-]{descent}
\externaldocument[perfect-]{perfect}
\externaldocument[more-morphisms-]{more-morphisms}
\externaldocument[flat-]{flat}
\externaldocument[groupoids-]{groupoids}
\externaldocument[more-groupoids-]{more-groupoids}
\externaldocument[etale-]{etale}
\externaldocument[chow-]{chow}
\externaldocument[intersection-]{intersection}
\externaldocument[pic-]{pic}
\externaldocument[weil-]{weil}
\externaldocument[adequate-]{adequate}
\externaldocument[dualizing-]{dualizing}
\externaldocument[duality-]{duality}
\externaldocument[discriminant-]{discriminant}
\externaldocument[derham-]{derham}
\externaldocument[local-cohomology-]{local-cohomology}
\externaldocument[algebraization-]{algebraization}
\externaldocument[curves-]{curves}
\externaldocument[resolve-]{resolve}
\externaldocument[models-]{models}
\externaldocument[functors-]{functors}
\externaldocument[equiv-]{equiv}
\externaldocument[pione-]{pione}
\externaldocument[etale-cohomology-]{etale-cohomology}
\externaldocument[proetale-]{proetale}
\externaldocument[relative-cycles-]{relative-cycles}
\externaldocument[more-etale-]{more-etale}
\externaldocument[trace-]{trace}
\externaldocument[crystalline-]{crystalline}
\externaldocument[spaces-]{spaces}
\externaldocument[spaces-properties-]{spaces-properties}
\externaldocument[spaces-morphisms-]{spaces-morphisms}
\externaldocument[decent-spaces-]{decent-spaces}
\externaldocument[spaces-cohomology-]{spaces-cohomology}
\externaldocument[spaces-limits-]{spaces-limits}
\externaldocument[spaces-divisors-]{spaces-divisors}
\externaldocument[spaces-over-fields-]{spaces-over-fields}
\externaldocument[spaces-topologies-]{spaces-topologies}
\externaldocument[spaces-descent-]{spaces-descent}
\externaldocument[spaces-perfect-]{spaces-perfect}
\externaldocument[spaces-more-morphisms-]{spaces-more-morphisms}
\externaldocument[spaces-flat-]{spaces-flat}
\externaldocument[spaces-groupoids-]{spaces-groupoids}
\externaldocument[spaces-more-groupoids-]{spaces-more-groupoids}
\externaldocument[bootstrap-]{bootstrap}
\externaldocument[spaces-pushouts-]{spaces-pushouts}
\externaldocument[spaces-chow-]{spaces-chow}
\externaldocument[groupoids-quotients-]{groupoids-quotients}
\externaldocument[spaces-more-cohomology-]{spaces-more-cohomology}
\externaldocument[spaces-simplicial-]{spaces-simplicial}
\externaldocument[spaces-duality-]{spaces-duality}
\externaldocument[formal-spaces-]{formal-spaces}
\externaldocument[restricted-]{restricted}
\externaldocument[spaces-resolve-]{spaces-resolve}
\externaldocument[formal-defos-]{formal-defos}
\externaldocument[defos-]{defos}
\externaldocument[cotangent-]{cotangent}
\externaldocument[examples-defos-]{examples-defos}
\externaldocument[algebraic-]{algebraic}
\externaldocument[examples-stacks-]{examples-stacks}
\externaldocument[stacks-sheaves-]{stacks-sheaves}
\externaldocument[criteria-]{criteria}
\externaldocument[artin-]{artin}
\externaldocument[quot-]{quot}
\externaldocument[stacks-properties-]{stacks-properties}
\externaldocument[stacks-morphisms-]{stacks-morphisms}
\externaldocument[stacks-limits-]{stacks-limits}
\externaldocument[stacks-cohomology-]{stacks-cohomology}
\externaldocument[stacks-perfect-]{stacks-perfect}
\externaldocument[stacks-introduction-]{stacks-introduction}
\externaldocument[stacks-more-morphisms-]{stacks-more-morphisms}
\externaldocument[stacks-geometry-]{stacks-geometry}
\externaldocument[moduli-]{moduli}
\externaldocument[moduli-curves-]{moduli-curves}
\externaldocument[examples-]{examples}
\externaldocument[exercises-]{exercises}
\externaldocument[guide-]{guide}
\externaldocument[desirables-]{desirables}
\externaldocument[coding-]{coding}
\externaldocument[obsolete-]{obsolete}
\externaldocument[fdl-]{fdl}
\externaldocument[index-]{index}

% Theorem environments.
%
\theoremstyle{plain}
\newtheorem{theorem}[subsection]{Theorem}
\newtheorem{proposition}[subsection]{Proposition}
\newtheorem{lemma}[subsection]{Lemma}

\theoremstyle{definition}
\newtheorem{definition}[subsection]{Definition}
\newtheorem{example}[subsection]{Example}
\newtheorem{exercise}[subsection]{Exercise}
\newtheorem{situation}[subsection]{Situation}

\theoremstyle{remark}
\newtheorem{remark}[subsection]{Remark}
\newtheorem{remarks}[subsection]{Remarks}

\numberwithin{equation}{subsection}

% Macros
%
\def\lim{\mathop{\mathrm{lim}}\nolimits}
\def\colim{\mathop{\mathrm{colim}}\nolimits}
\def\Spec{\mathop{\mathrm{Spec}}}
\def\Hom{\mathop{\mathrm{Hom}}\nolimits}
\def\Ext{\mathop{\mathrm{Ext}}\nolimits}
\def\SheafHom{\mathop{\mathcal{H}\!\mathit{om}}\nolimits}
\def\SheafExt{\mathop{\mathcal{E}\!\mathit{xt}}\nolimits}
\def\Sch{\mathit{Sch}}
\def\Mor{\mathop{\mathrm{Mor}}\nolimits}
\def\Ob{\mathop{\mathrm{Ob}}\nolimits}
\def\Sh{\mathop{\mathit{Sh}}\nolimits}
\def\NL{\mathop{N\!L}\nolimits}
\def\CH{\mathop{\mathrm{CH}}\nolimits}
\def\proetale{{pro\text{-}\acute{e}tale}}
\def\etale{{\acute{e}tale}}
\def\QCoh{\mathit{QCoh}}
\def\Ker{\mathop{\mathrm{Ker}}}
\def\Im{\mathop{\mathrm{Im}}}
\def\Coker{\mathop{\mathrm{Coker}}}
\def\Coim{\mathop{\mathrm{Coim}}}

% Boxtimes
%
\DeclareMathSymbol{\boxtimes}{\mathbin}{AMSa}{"02}

%
% Macros for moduli stacks/spaces
%
\def\QCohstack{\mathcal{QC}\!\mathit{oh}}
\def\Cohstack{\mathcal{C}\!\mathit{oh}}
\def\Spacesstack{\mathcal{S}\!\mathit{paces}}
\def\Quotfunctor{\mathrm{Quot}}
\def\Hilbfunctor{\mathrm{Hilb}}
\def\Curvesstack{\mathcal{C}\!\mathit{urves}}
\def\Polarizedstack{\mathcal{P}\!\mathit{olarized}}
\def\Complexesstack{\mathcal{C}\!\mathit{omplexes}}
% \Pic is the operator that assigns to X its picard group, usage \Pic(X)
% \Picardstack_{X/B} denotes the Picard stack of X over B
% \Picardfunctor_{X/B} denotes the Picard functor of X over B
\def\Pic{\mathop{\mathrm{Pic}}\nolimits}
\def\Picardstack{\mathcal{P}\!\mathit{ic}}
\def\Picardfunctor{\mathrm{Pic}}
\def\Deformationcategory{\mathcal{D}\!\mathit{ef}}
\usepackage{cleveref}
\crefname{originaldefinition}{original definition}{original definitions}
\Crefname{originaldefinition}{Original definition}{Original definitions}
\makeatletter
\newlabel{morphisms-definition-finite-type-point}{{02J1}{source}{}{sourcepointer-1}{}}
\newlabel{morphisms-definition-finite-type-point@cref}{{[originaldefinition][0][]02J1}{[1][0][]source}{}{}{}}
\makeatother
\crefname{originalreference}{original reference}{original references}
\Crefname{originalreference}{Original reference}{Original references}
\makeatletter
\newlabel{situation-slice}{{situation-slice}{source}{}{sourcepointer-2}{}}
\newlabel{situation-slice@cref}{{[originalreference][0][]situation-slice}{[1][0][]source}{}{}{}}
\makeatother
\crefname{originallemma}{original lemma}{original lemmas}
\Crefname{originallemma}{Original lemma}{Original lemmas}
\makeatletter
\newlabel{morphisms-lemma-identify-finite-type-points}{{02J2}{source}{}{sourcepointer-3}{}}
\newlabel{morphisms-lemma-identify-finite-type-points@cref}{{[originallemma][0][]02J2}{[1][0][]source}{}{}{}}
\makeatother
\crefname{originallemma}{original lemma}{original lemmas}
\Crefname{originallemma}{Original lemma}{Original lemmas}
\makeatletter
\newlabel{morphisms-lemma-dimension-fibre-at-a-point}{{02FX}{source}{}{sourcepointer-4}{}}
\newlabel{morphisms-lemma-dimension-fibre-at-a-point@cref}{{[originallemma][0][]02FX}{[1][0][]source}{}{}{}}
\makeatother
\crefname{originallemma}{original lemma}{original lemmas}
\Crefname{originallemma}{Original lemma}{Original lemmas}
\makeatletter
\newlabel{lemma-restrict-preserves-type}{{lemma-restrict-preserves-type}{source}{}{sourcepointer-5}{}}
\newlabel{lemma-restrict-preserves-type@cref}{{[originallemma][0][]lemma-restrict-preserves-type}{[1][0][]source}{}{}{}}
\makeatother
\crefname{originallemma}{original lemma}{original lemmas}
\Crefname{originallemma}{Original lemma}{Original lemmas}
\makeatletter
\newlabel{lemma-groupoid-on-field-dimension-equal-stabilizer}{{lemma-groupoid-on-field-dimension-equal-stabilizer}{source}{}{sourcepointer-6}{}}
\newlabel{lemma-groupoid-on-field-dimension-equal-stabilizer@cref}{{[originallemma][0][]lemma-groupoid-on-field-dimension-equal-stabilizer}{[1][0][]source}{}{}{}}
\makeatother
\crefname{originallemma}{original lemma}{original lemmas}
\Crefname{originallemma}{Original lemma}{Original lemmas}
\makeatletter
\newlabel{lemma-groupoid-on-field-locally-finite-type-dimension}{{lemma-groupoid-on-field-locally-finite-type-dimension}{source}{}{sourcepointer-7}{}}
\newlabel{lemma-groupoid-on-field-locally-finite-type-dimension@cref}{{[originallemma][0][]lemma-groupoid-on-field-locally-finite-type-dimension}{[1][0][]source}{}{}{}}
\makeatother
\crefname{originallemma}{original lemma}{original lemmas}
\Crefname{originallemma}{Original lemma}{Original lemmas}
\makeatletter
\newlabel{groupoids-lemma-restrict-stabilizer}{{04ML}{source}{}{sourcepointer-8}{}}
\newlabel{groupoids-lemma-restrict-stabilizer@cref}{{[originallemma][0][]04ML}{[1][0][]source}{}{}{}}
\makeatother
\crefname{originallemma}{original lemma}{original lemmas}
\Crefname{originallemma}{Original lemma}{Original lemmas}
\makeatletter
\newlabel{algebra-lemma-CM-dim-formula}{{00NA}{source}{}{sourcepointer-9}{}}
\newlabel{algebra-lemma-CM-dim-formula@cref}{{[originallemma][0][]00NA}{[1][0][]source}{}{}{}}
\makeatother
\crefname{originallemma}{original lemma}{original lemmas}
\Crefname{originallemma}{Original lemma}{Original lemmas}
\makeatletter
\newlabel{algebra-lemma-silly}{{00DS}{source}{}{sourcepointer-10}{}}
\newlabel{algebra-lemma-silly@cref}{{[originallemma][0][]00DS}{[1][0][]source}{}{}{}}
\makeatother
\crefname{originallemma}{original lemma}{original lemmas}
\Crefname{originallemma}{Original lemma}{Original lemmas}
\makeatletter
\newlabel{algebra-lemma-one-equation}{{00KW}{source}{}{sourcepointer-11}{}}
\newlabel{algebra-lemma-one-equation@cref}{{[originallemma][0][]00KW}{[1][0][]source}{}{}{}}
\makeatother
\crefname{originallemma}{original lemma}{original lemmas}
\Crefname{originallemma}{Original lemma}{Original lemmas}
\makeatletter
\newlabel{algebra-lemma-reformulate-CM}{{02JN}{source}{}{sourcepointer-12}{}}
\newlabel{algebra-lemma-reformulate-CM@cref}{{[originallemma][0][]02JN}{[1][0][]source}{}{}{}}
\makeatother
\crefname{originallemma}{original lemma}{original lemmas}
\Crefname{originallemma}{Original lemma}{Original lemmas}
\makeatletter
\newlabel{algebra-lemma-grothendieck-general}{{046Z}{source}{}{sourcepointer-13}{}}
\newlabel{algebra-lemma-grothendieck-general@cref}{{[originallemma][0][]046Z}{[1][0][]source}{}{}{}}
\makeatother
\crefname{originalequation}{original equation}{original equations}
\Crefname{originalequation}{Original equation}{Original equations}
\makeatletter
\newlabel{equation-restriction}{{equation-restriction}{source}{}{sourcepointer-14}{}}
\newlabel{equation-restriction@cref}{{[originalequation][0][]equation-restriction}{[1][0][]source}{}{}{}}
\makeatother
\title{Slicing: an lfp immersion cutting the source fibre dimension by one when the stabilizer is smaller}
\author{The Stacks Project Authors}
\date{stacks-project:more-groupoids.tex:local-snapshot-20261012; GNU FDL 1.2 or later}
\begin{document}
\maketitle
\setcounter{subsection}{1}
\noindent\textit{Editorial note.} All delivered mathematical text is the author's own, copied from the pinned Stacks Project chapter file more-groupoids.tex; no statement or proof step was rewritten, reordered or removed. The counted span is the author environment and its own proof environment. References inside the span are handed to the assembler as navigation stubs that carry the canonical tag of the cited tag only; whether each of them is load-bearing is recorded on the stub itself.\par\medskip
\noindent\textit{Allowed standard prerequisites (editorial).}\begin{itemize}
\item Definition of a finite type point (morphisms chapter, Definition morphisms-definition-finite-type-point, canonical tag 02J1). cited at statement lines 2072-2073. The hypothesis on u is phrased with this notion.
\item Notation situation for slicing: the restriction of the groupoid via g, the stabilizer group scheme G, the map h equal to s composed with the first projection, and the fibres F\_u, G\_u, F'\_\{u'\} (more-groupoids chapter, Situation situation-slice, no canonical tag in the snapshot; no-tag pointer to the source label line). cited at statement lines 2075 and 2083. Fixes the notation d\_1, d\_2, h, F\_u, G\_u and the conclusion items in which it is phrased.
\item A finite type point has an affine neighbourhood in which it is closed (morphisms chapter, Lemma morphisms-lemma-identify-finite-type-points, canonical tag 02J2). cited at proof lines 2095-2097. Provides the affine neighbourhood Spec(A) with u a closed point.
\item Dimension of the fibre of a morphism at a point equals the dimension of the local ring of the fibre (morphisms chapter, Lemma morphisms-lemma-dimension-fibre-at-a-point, canonical tag 02FX). cited at proof lines 2128-2129 and 2229-2232. Identifies d\_2 with the dimension of the local ring of the fibre F\_u at e(u), and later evaluates the dimension of F'\_u at e'(u) as the dimension of the cut fibre ring.
\item Restriction of a groupoid preserves local finite type of the structure maps (more-groupoids chapter, Lemma lemma-restrict-preserves-type, no canonical tag in the snapshot; no-tag pointer to the source label line). cited at proof lines 2134-2137. The restriction of R to the point u is a groupoid scheme with locally finite type structure maps.
\item For a groupoid on a field with locally finite type structure maps the dimensions of the restricted arrow scheme and the stabilizer agree (more-groupoids chapter, Lemma lemma-groupoid-on-field-dimension-equal-stabilizer, no canonical tag in the snapshot; no-tag pointer to the source label line). cited at proof lines 2138-2140. Gives dim(R'') = dim(G'') for the restriction to the point u.
\item For such groupoids on a field the dimension equals the dimension of the local ring at the identity point (more-groupoids chapter, Lemma lemma-groupoid-on-field-locally-finite-type-dimension, no canonical tag in the snapshot; no-tag pointer to the source label line). cited at proof lines 2141-2142. Converts dim(R'') into the dimension of the local ring of R'' at e''.
\item The stabilizer of a restriction is the fibre of the stabilizer (groupoids chapter, Lemma groupoids-lemma-restrict-stabilizer, canonical tag 04ML). cited at proof lines 2143-2145. Identifies G'' with G\_u, so the dimension of the local ring of R'' at e'' equals d\_1.
\item Dimension formula for Cohen-Macaulay local rings and minimal primes (algebra chapter, Lemma algebra-lemma-CM-dim-formula, canonical tag 00NA). cited at proof lines 2181-2183. Each minimal prime of the Cohen-Macaulay fibre ring has quotient of dimension d\_2.
\item Prime avoidance (algebra chapter, Lemma algebra-lemma-silly, canonical tag 00DS). cited at proof lines 2192-2194. Produces f in the maximal ideal whose t-image avoids all the minimal primes.
\item One equation cuts dimension (algebra chapter, Lemma algebra-lemma-one-equation, canonical tag 00KW). cited at proof lines 2195-2197. Gives the displayed inequality dim below d\_2 for the fibre ring cut by f.
\item Reformulation of Cohen-Macaulayness: over a Cohen-Macaulay ring a parameter element is a non-zerodivisor and cutting by it stays Cohen-Macaulay of dimension one less (algebra chapter, Lemma algebra-lemma-reformulate-CM, canonical tag 02JN). cited at proof lines 2213-2217, 2223-2224 and 2233-2234. t(f) is a non-zerodivisor on the Cohen-Macaulay fibre ring, the cut ring is Cohen-Macaulay, and its dimension is d\_2 minus one.
\item Grothendieck flatness criterion by fibers (algebra chapter, Lemma algebra-lemma-grothendieck-general, canonical tag 046Z). cited at proof lines 2218-2220. Concludes flatness of the local ring map with Cohen-Macaulay fibre ring, giving part (4).
\item Defining diagram of a restriction of a groupoid (more-groupoids chapter, Equation equation-restriction, no canonical tag in the snapshot; no-tag pointer to the source label line). cited at proof line 2226. Identifies the fibre F'\_u with the fibre of h over u, reducing part (5) to the computed dimension.
\end{itemize}\medskip
\ifdefined\hypertarget\hypertarget{sourcepointer-1}{}\fi
\noindent\textit{Source reference.} Original definition 02J1: sources/stacks-chow/tags.txt, source lines 3320--3320. Cited at statement lines 2072-2073: the hypothesis on u is phrased with this notion.\par\medskip
\ifdefined\hypertarget\hypertarget{sourcepointer-2}{}\fi
\noindent\textit{Source reference.} Original reference situation-slice: sources/stacks-more-groupoids/pinned-original-source.tex, source lines 2041--2041. Cited at statement lines 2075 and 2083: fixes the notation d\_1, d\_2, h and the conclusion items.\par\medskip
\ifdefined\hypertarget\hypertarget{sourcepointer-3}{}\fi
\noindent\textit{Source reference.} Original lemma 02J2: sources/stacks-chow/tags.txt, source lines 3321--3321. Cited at proof lines 2095-2097: provides the affine neighbourhood with u closed.\par\medskip
\ifdefined\hypertarget\hypertarget{sourcepointer-4}{}\fi
\noindent\textit{Source reference.} Original lemma 02FX: sources/stacks-chow/tags.txt, source lines 3208--3208. Cited at proof lines 2128-2129 and 2229-2232: identifies d\_2 with the local fibre ring dimension and evaluates the dimension of F'\_u at e'(u).\par\medskip
\ifdefined\hypertarget\hypertarget{sourcepointer-5}{}\fi
\noindent\textit{Source reference.} Original lemma lemma-restrict-preserves-type: sources/stacks-more-groupoids/pinned-original-source.tex, source lines 923--923. Cited at proof lines 2134-2137: the restriction to the point u has locally finite type structure maps.\par\medskip
\ifdefined\hypertarget\hypertarget{sourcepointer-6}{}\fi
\noindent\textit{Source reference.} Original lemma lemma-groupoid-on-field-dimension-equal-stabilizer: sources/stacks-more-groupoids/pinned-original-source.tex, source lines 1554--1554. Cited at proof lines 2138-2140: dim of the restricted arrow scheme equals dim of its stabilizer.\par\medskip
\ifdefined\hypertarget\hypertarget{sourcepointer-7}{}\fi
\noindent\textit{Source reference.} Original lemma lemma-groupoid-on-field-locally-finite-type-dimension: sources/stacks-more-groupoids/pinned-original-source.tex, source lines 1515--1515. Cited at proof lines 2141-2142: converts the dimension into the local ring dimension at the identity.\par\medskip
\ifdefined\hypertarget\hypertarget{sourcepointer-8}{}\fi
\noindent\textit{Source reference.} Original lemma 04ML: sources/stacks-chow/tags.txt, source lines 6047--6047. Cited at proof lines 2143-2145: identifies G'' with G\_u, giving dimension d\_1.\par\medskip
\ifdefined\hypertarget\hypertarget{sourcepointer-9}{}\fi
\noindent\textit{Source reference.} Original lemma 00NA: sources/stacks-chow/tags.txt, source lines 852--852. Cited at proof lines 2181-2183: each minimal prime quotient of the CM fibre ring has dimension d\_2.\par\medskip
\ifdefined\hypertarget\hypertarget{sourcepointer-10}{}\fi
\noindent\textit{Source reference.} Original lemma 00DS: sources/stacks-chow/tags.txt, source lines 508--508. Cited at proof lines 2192-2194: produces f in the maximal ideal avoiding all minimal primes after t.\par\medskip
\ifdefined\hypertarget\hypertarget{sourcepointer-11}{}\fi
\noindent\textit{Source reference.} Original lemma 00KW: sources/stacks-chow/tags.txt, source lines 766--766. Cited at proof lines 2195-2197: gives the displayed inequality below d\_2.\par\medskip
\ifdefined\hypertarget\hypertarget{sourcepointer-12}{}\fi
\noindent\textit{Source reference.} Original lemma 02JN: sources/stacks-chow/tags.txt, source lines 3342--3342. Cited at proof lines 2213-2217, 2223-2224 and 2233-2234: non-zerodivisibility of t(f), CM-ness of the cut ring, and the dimension drop by one.\par\medskip
\ifdefined\hypertarget\hypertarget{sourcepointer-13}{}\fi
\noindent\textit{Source reference.} Original lemma 046Z: sources/stacks-chow/tags.txt, source lines 5490--5490. Cited at proof lines 2218-2220: flatness of the local map with Cohen-Macaulay fibre ring, giving part (4).\par\medskip
\ifdefined\hypertarget\hypertarget{sourcepointer-14}{}\fi
\noindent\textit{Source reference.} Original equation equation-restriction: sources/stacks-more-groupoids/pinned-original-source.tex, source lines 905--905. Cited at proof line 2226: identifies the fibre F'\_u with the fibre of h over u.\par\medskip
\begin{lemma}
\label{lemma-slice}
Let $S$ be a scheme.
Let $(U, R, s, t, c, e, i)$ be a groupoid scheme over $S$.
Let $G \to U$ be the stabilizer group scheme.
Assume $s$ and $t$ are Cohen-Macaulay and locally of finite presentation.
Let $u \in U$ be a finite type point of the scheme $U$, see
Morphisms, Definition \ref{morphisms-definition-finite-type-point}.
With notation as in
Situation \ref{situation-slice},
set
$$
d_1 = \dim(G_u), \quad
d_2 = \dim_{e(u)}(F_u).
$$
If $d_2 > d_1$, then there exist an affine scheme $U'$
and a morphism $g : U' \to U$ such that (with notation as in
Situation \ref{situation-slice})
\begin{enumerate}
\item $g$ is an immersion
\item $u \in U'$,
\item $g$ is locally of finite presentation,
\item the morphism $h : U' \times_{g, U, t} R \longrightarrow U$
is Cohen-Macaulay at $(u, e(u))$, and
\item we have $\dim_{e'(u)}(F'_u) = d_2 - 1$.
\end{enumerate}
\end{lemma}

\begin{proof}
Let $\Spec(A) \subset U$ be an affine neighbourhood of $u$
such that $u$ corresponds to a closed point of $U$, see
Morphisms, Lemma \ref{morphisms-lemma-identify-finite-type-points}.
Let $\Spec(B) \subset R$ be an affine neighbourhood of $e(u)$
which maps via $j$ into the open
$\Spec(A) \times_S \Spec(A) \subset U \times_S U$.
Let $\mathfrak m \subset A$ be the maximal ideal corresponding to $u$.
Let $\mathfrak q \subset B$ be the prime ideal corresponding to $e(u)$.
Pictures:
$$
\vcenter{
\xymatrix{
B & A \ar[l]^s \\
A \ar[u]^t
}
}
\quad\text{and}\quad
\vcenter{
\xymatrix{
B_{\mathfrak q} & A_{\mathfrak m} \ar[l]^s \\
A_{\mathfrak m} \ar[u]^t
}
}
$$
Note that the two induced maps
$s, t : \kappa(\mathfrak m) \to \kappa(\mathfrak q)$
are equal and isomorphisms as $s \circ e = t \circ e = \text{id}_U$.
In particular we see that $\mathfrak q$
is a maximal ideal as well. The ring maps $s, t : A \to B$ are
of finite presentation and flat. By assumption the ring
$$
\mathcal{O}_{F_u, e(u)} = B_{\mathfrak q}/s(\mathfrak m)B_{\mathfrak q}
$$
is Cohen-Macaulay of dimension $d_2$. The equality of dimension holds by
Morphisms, Lemma \ref{morphisms-lemma-dimension-fibre-at-a-point}.

\medskip\noindent
Let $R''$ be the restriction of $R$ to $u = \Spec(\kappa(u))$
via the morphism $\Spec(\kappa(u)) \to U$.
As $u \to U$ is locally of finite type,
we see that $(\Spec(\kappa(u)), R'', s'', t'', c'')$
is a groupoid scheme with $s'', t''$ locally of finite type, see
Lemma \ref{lemma-restrict-preserves-type}.
By
Lemma \ref{lemma-groupoid-on-field-dimension-equal-stabilizer}
this implies that $\dim(G'') = \dim(R'')$. We also have
$\dim(R'') = \dim_{e''}(R'') = \dim(\mathcal{O}_{R'', e''})$, see
Lemma \ref{lemma-groupoid-on-field-locally-finite-type-dimension}.
By
Groupoids, Lemma \ref{groupoids-lemma-restrict-stabilizer}
we have $G'' = G_u$. Hence we conclude that
$\dim(\mathcal{O}_{R'', e''}) = d_1$.

\medskip\noindent
As a scheme $R''$ is
$$
R'' =
R \times_{(U \times_S U)}
\Big(
\Spec(\kappa(\mathfrak m)) \times_S \Spec(\kappa(\mathfrak m))
\Big)
$$
Hence an affine open neighbourhood of $e''$ is the spectrum of the ring
$$
B \otimes_{(A \otimes A)} (\kappa(\mathfrak m) \otimes \kappa(\mathfrak m))
=
B/s(\mathfrak m)B + t(\mathfrak m)B
$$
We conclude that
$$
\mathcal{O}_{R'', e''} =
B_{\mathfrak q}/s(\mathfrak m)B_{\mathfrak q} + t(\mathfrak m)B_{\mathfrak q}
$$
and so now we know that this ring has dimension $d_1$.

\medskip\noindent
We claim this implies we can find
an element $f \in \mathfrak m$ such that
$$
\dim(B_{\mathfrak q}/(s(\mathfrak m)B_{\mathfrak q} + fB_{\mathfrak q}) < d_2
$$
Namely, suppose $\mathfrak n_j \supset s(\mathfrak m)B_{\mathfrak q}$,
$j = 1, \ldots, m$ correspond to the minimal primes of the local ring
$B_{\mathfrak q}/s(\mathfrak m)B_{\mathfrak q}$.
There are finitely many as this ring is Noetherian (since it is essentially
of finite type over a field -- but also because a Cohen-Macaulay ring is
Noetherian). By the Cohen-Macaulay condition we have
$\dim(B_{\mathfrak q}/\mathfrak n_j) = d_2$, for example by
Algebra, Lemma \ref{algebra-lemma-CM-dim-formula}.
Note that
$\dim(B_{\mathfrak q}/(\mathfrak n_j + t(\mathfrak m)B_{\mathfrak q}))
\leq d_1$
as it is a quotient of the ring
$\mathcal{O}_{R'', e''} =
B_{\mathfrak q}/s(\mathfrak m)B_{\mathfrak q} + t(\mathfrak m)B_{\mathfrak q}$
which has dimension $d_1$. As $d_1 < d_2$ this implies that
$\mathfrak m \not \subset t^{-1}(\mathfrak n_i)$.
By prime avoidance, see
Algebra, Lemma \ref{algebra-lemma-silly},
we can find $f \in \mathfrak m$ with $t(f) \not \in \mathfrak n_j$ for
$j = 1, \ldots, m$. For this choice of $f$ we have
the displayed inequality above, see
Algebra, Lemma \ref{algebra-lemma-one-equation}.

\medskip\noindent
Set $A' = A/fA$ and $U' = \Spec(A')$. Then it is clear that
$U' \to U$ is an immersion, locally of finite presentation
and that $u \in U'$. Thus (1), (2) and (3) of the lemma hold.
The morphism
$$
U' \times_{g, U, t} R \longrightarrow U
$$
factors through $\Spec(A)$ and corresponds to the ring map
$$
\xymatrix{
B/t(f)B \ar@{=}[r] & A/(f) \otimes_{A, t} B & A \ar[l]_-s
}
$$
Now, we see $t(f)$ is not a zerodivisor on
$B_{\mathfrak q}/s(\mathfrak m)B_{\mathfrak q}$ as this is a
Cohen-Macaulay ring of positive dimension and $f$ is not contained
in any minimal prime, see for example
Algebra, Lemma \ref{algebra-lemma-reformulate-CM}.
Hence by
Algebra, Lemma \ref{algebra-lemma-grothendieck-general}
we conclude that $s : A_{\mathfrak m} \to B_{\mathfrak q}/t(f)B_{\mathfrak q}$
is flat with fibre ring
$B_{\mathfrak q}/(s(\mathfrak m)B_{\mathfrak q} + t(f)B_{\mathfrak q})$
which is Cohen-Macaulay by
Algebra, Lemma \ref{algebra-lemma-reformulate-CM}
again. This implies part (4) of the lemma.
To see part (5) note that by Diagram (\ref{equation-restriction})
the fibre $F'_u$ is equal to the fibre of $h$ over $u$.
Hence
$\dim_{e'(u)}(F'_u) =
\dim(B_{\mathfrak q}/(s(\mathfrak m)B_{\mathfrak q} + t(f)B_{\mathfrak q}))$
by
Morphisms, Lemma \ref{morphisms-lemma-dimension-fibre-at-a-point}
and the dimension of this ring is $d_2 - 1$ by
Algebra, Lemma \ref{algebra-lemma-reformulate-CM}
once more. This proves the final assertion of the lemma and we win.
\end{proof}

\end{document}
